\subsection{整数与有限集上的运算}

命题1. 设 \((a_i)_{i\in I}\) 是一个整数的有限族。则基数 \(\sum_{i\in I}a_i\) 与 \(\prod_{i\in I}a_i\) 是整数。

让我们首先证明，如果 \(a\) 与 \(b\) 是整数，则 \(a + b\) 也是整数。我们对 \(b\) 进行归纳。该命题在 \(b = 0\) 时成立，因为 \(a + 0 = a\) 。如果 \(a + b\) 是整数，则 \((a + b) + 1\) 也是整数 ( \(\S 4\) ，第1号，命题1）。但 \((a + b) + 1 = a + (b + 1)\) ( \(\S 3\) ，第3号，命题5的推论）；因此 \(a + (b + 1)\) 是整数，从而 \(a + b\) 对所有整数 \(b\) 而言都是整数.

接下来对 \(n = \operatorname{Card}(\mathbf{I})\) 作归纳，证明 \(\sum_{i \in \mathbf{I}} a_i\) 是整数。当 \(n = 0\) 时这是显然的，因为 \(\mathbf{I} = \emptyset\) 且 \(\sum_{i \in \mathbf{I}} a_i = 0\) 。如果 \(\operatorname{Card}(\mathbf{I}) = n + 1\) ，可写成 \(\mathbf{I} = \mathbf{J} \cup \{k\}\) ，其中 \(\operatorname{Card}(\mathbf{J}) = n\) 且 \(k \notin \mathbf{J}\) 。于是

\[
\sum_ {i \in \mathbf {I}} a _ {i} = a _ {k} + \sum_ {i \in \mathbf {J}} a _ {i}
\]

（§ 3，第 3 号，命题 5）。归纳假设是 \(\sum_{i\in J}a_i\) 是整数；因此 \(a_{k} + \sum_{i\in J}a_{i}\) 也是整数，这由证明的第一段得出。这表明 \(\sum_{i\in I}a_{i}\) 对所有 \(n\) 都是整数.

由于两个整数a和b的乘积ab是有限个全等于a的整数之和（§3，第4号，命题6，推论2），因此ab是整数。我们将通过对n = Card (I)进行归纳来证明 \(\prod_{i\in I}a_{i}\) 是一个整数。这对n = 0成立，因为此时

\[
\prod_ {i \in \mathbf {I}} a _ {i} = 1.
\]

若 Card (I) = n + 1，则（沿用上述记号）我们有

\[
\prod_ {i \in \mathbf {I}} a _ {i} = a _ {k}. \prod_ {i \in \mathbf {J}} a _ {i}
\]

（§3，第3号，命题5），因此归纳假设蕴含 \(\prod_{i\in I}a_i\) 是一个整数。因此 \(\prod_{i\in I}a_i\) 对所有 \(n\) 都是整数.

推论1. 有限集合的有限族 \((\mathbf{X}_i)_{i\in \mathbf{I}}\) 的并集E是有限集合。

因为族 \((\mathbf{X}_{i})\) 的和 S 是有限的；并且由于存在从 S 到 E 的满射（第二章，§4，第 8 号），集合 E 是有限的（§4，第 2 号，命题 2，推论 3）。

推论 2. 有限个有限集合的乘积是有限集合。

推论3。如果 \(a\) 与 \(b\) 是整数，则 \(a^b\) 是一个整数。

因为 \(a^b\) 是一个有限整数族的乘积，该族中每个整数都等于 \(a\) （§3，第5号，命题10）。

推论4. 有限集合E的子集所构成的集合是有限的。

因为其基数为 \(2^{\text{Card (E)}}\) ( \(\S 3\) ，第5号，命题12）。

